\chapter{Why Water Is So Special}

\begin{kaichang}
Pour a glass of ice water and add three ice cubes. Watch for a minute, and you will see three strange things.

First, ice floats. Consider how odd that is: a substance's solid is lighter than its own liquid. Solid wax sinks in melted wax; solid olive oil sinks in liquid olive oil. Almost every substance has a denser solid that sinks. Water does the opposite.

Second, droplets slowly appear on the glass's outside. It is not leaking, so where did they come from?

Third, look again half an hour later. The ice has melted, yet the water level is almost unchanged. The bit of ice above the surface seems to have ``vanished'' on melting.

Record these three oddities and add something you may not have observed but can guess: why do adults sponge you with warm water when you have a fever? Why are coastal summers cooler and winters warmer than inland ones? The answers all hide in a molecule containing only three atoms.
\end{kaichang}

\section{An ``Anomalous Substance's'' Charge Sheet}

In physics and chemistry classrooms, water is usually a ``standard substance'': a unit for specific heat, a density reference, and a benchmark of acid--base neutrality. Yet scientists see it as \textbf{the most anomalous common substance}. Let us list its ``offenses'':

\begin{itemize}[itemsep=2pt]
\item \textbf{It is far too fond of being liquid}. Molecules of similar ``weight''---methane (\chem{CH_4}), ammonia (\chem{NH_3}), and hydrogen sulfide (\chem{H_2S})---are gases at room temperature. Judging by its relatives, water should boil around a hundred degrees below zero, leaving Earth without rivers, lakes, or oceans. Yet it boils at $100\,^{\circ}\mathrm{C}$.
\item \textbf{Its solid is lighter than its liquid}. Ice's density is about \ud{0.917}{g/cm^3}, so it floats with roughly a tenth of its volume above water: literally the ``tip of the iceberg.''
\item \textbf{It takes heat remarkably well}. Heat equal kettles of water and oil over identical flames, and oil warms much faster. Water absorbs almost more heat per $1\,^{\circ}\mathrm{C}$ rise than any other common liquid.
\item \textbf{It is picky}. Salt, sugar, and alcohol vanish with a stir, but oil insists on separating, leaving salad dressing in two layers after standing.
\end{itemize}

Is one culprit behind these offenses? Yes. To catch it, zoom into a single water molecule. These four apparently unrelated quirks are the same structural cause appearing in different settings.

\section{A Bent Molecule with Unequal Ends}

Chapter 19 explained oxygen--hydrogen polar covalent bonds: electronegative oxygen pulls shared electrons closer, making oxygen slightly negative and hydrogen slightly positive. Chapter 20 showed that water is bent, not linear, with its hydrogens about $104.5^{\circ}$ apart like Mickey Mouse's ears.

Together, these facts have major consequences: the positive and negative charge centers \textbf{cannot coincide}. If water were linear and symmetric like carbon dioxide, its partial charges would cancel externally. But it bends, leaving ``Mickey's chin'' (oxygen) negative and his ``ear tips'' (hydrogen) positive. A water molecule resembles a tiny crooked magnet, negative at one end and positive at the other.

Chapter 22 introduced what follows: one molecule's positive ``ear tip'' attracts a neighbor's negative ``chin,'' forming a \textbf{hydrogen bond}, weaker than a chemical bond but much stronger than ordinary intermolecular forces. Unlike an internal electron-sharing covalent bond, this is electrostatic attraction \textbf{between} molecules, about one-twentieth as strong.

The key is how many ``hands'' water has. Oxygen retains two nonbonding lone pairs that can accept hydrogens from two other molecules, while its own two hydrogens reach toward two more neighbors. Altogether, one molecule can hold \textbf{four} neighbors simultaneously. The four directions point toward tetrahedral vertices, no coincidence but the very reason those ``Mickey Mouse ears'' approach a tetrahedral angle.

Liquid water is therefore not a bowl of independent ``molecular beans,'' but a continually renewed net. Each molecule extends hydrogen bonds to its neighbors, bonds continually breaking and immediately reconnecting, lasting on average only a trillionth of a second. Yet at any instant most molecules remain attached. All water's quirks are projections of this net.

\section{One Network Explains Four Offenses}

\subsection{Offense One: An Absurdly High Boiling Point}

Boiling essentially gives molecules enough speed to escape neighbors and enter the air. Methane molecules have only weak dispersion (Chapter 22) and separate with a gentle push, boiling at $-161\,^{\circ}\mathrm{C}$. Escaping water must first tear away its hydrogen bonds, costing much more energy. Only at $100\,^{\circ}\mathrm{C}$ does widespread ``separation'' begin. Water without hydrogen bonds is a hypothetical substance impossible on Earth; with them, it remains liquid at ordinary temperatures, providing oceans, clouds, rain, and over half your body weight.

\subsection{Offense Two: Floating Ice}

Liquid water's network is disordered and constantly rearranging, allowing molecules to squeeze into gaps. Freezing changes this: weakened thermal motion no longer overcomes hydrogen bonds' directional requirements, and each molecule firmly holds four neighbors in an orderly tetrahedral scaffold. Accommodating the bond directions leaves large holes, like a sparse framework built to align its joints. Thus \textbf{ice has a more open structure than liquid water}. The same molecules occupy more space, lowering density and making ice float. Freezing expands water by about 9\%; that 9\% bursts pipes and fractures rocks in winter.

This is wonderful news for life. Winter ice forms on lakes' surfaces and floats like an insulating blanket, keeping water liquid below so fish and aquatic plants survive. If ice sank like other solids, lakes would freeze upward from the bottom, radically changing northern aquatic ecosystems.

\subsection{Offense Three: Taking Heat in Stride}

Where does heat added to water go? In ordinary liquids, most directly increases random molecular motion, quickly raising temperature. In water, incoming energy must first do ``housework'': breaking hydrogen bonds and loosening the net. Only afterward does the remainder raise temperature. Hence water's large heat capacity: raising one gram by $1\,^{\circ}\mathrm{C}$ requires about \ud{4.2}{J}, roughly twice ethanol's requirement and nine times iron's.

Conversely, evaporation lets the fastest molecules escape the whole net, carrying substantial energy away: each gram removes about \ud{2400}{J}. That is why sweating and warm-water sponging reduce fever: evaporation ``express-ships'' body heat into the air.

This combination of heat capacity and evaporative cooling is also your temperature-regulation system. About 60 percent of your body is water. Your muscular and internal-organ ``engine'' continually produces heat, while body fluids' large heat capacity acts as a buffer reservoir, preventing meals or runs from dramatically changing core temperature. When the reservoir nears capacity, sweat glands start evaporation and send excess heat outside. Mammals' ability to maintain temperatures near $37\,^{\circ}\mathrm{C}$ year-round rests first on being filled with this heat-absorbing liquid.

\begin{qiaoliang}
Calculate for one glass. Heating about 200 grams of water from $20\,^{\circ}\mathrm{C}$ to $100\,^{\circ}\mathrm{C}$, a rise of $80\,^{\circ}\mathrm{C}$, requires approximately
\[200 \times 4.2 \times 80 \approx 6.7\times 10^{4}\ \mathrm{J}.\]
Vaporizing that same glass entirely from $100\,^{\circ}\mathrm{C}$ requires about $200 \times 2250 \approx 4.5\times 10^{5}\ \mathrm{J}$, seven times as much. You have seen this: steady heat brings water to a rolling boil within minutes, but drying it out takes far longer. Most energy goes to ``dismantling the net,'' releasing molecules from hydrogen bonds, rather than ``speeding them up'' and raising temperature. Coastal mildness has the same account: seawater absorbs much heat by day and in summer with little warming, then releases it slowly at night and in winter. The ocean is a giant climate thermos.
\end{qiaoliang}

\subsection{Offense Four: Pickiness}

Salt disappears in water while oil gathers into layers. Both follow the net's ``friendship rules.''

Salt first: its sodium--chloride crystal (Chapter 18) meets water's negative ends surrounding sodium and positive ends surrounding chloride. Groups of water molecules ``lift'' an ion off the lattice and wrap it in a little watery ball: \textbf{hydration}. Ion--water attraction compensates for lattice separation and network rearrangement, so salt dissolves. Sugar has no charge, but its many water-loving hydroxyl groups can join the hydrogen-bond net directly and are equally welcome. The rule is \textbf{like dissolves like}: polar water readily accepts charged or polar partners.

This is useful daily. Engine oil and paint on your hands resist water but wipe off with gasoline or alcohol, because nonpolar hydrocarbon dirt needs similarly nonpolar solvents. Dry cleaners likewise replace water with nonpolar solvents to wash whole garments. Cooking is similar: salt, sugar, and vinegar dissolve in soup because they like water, while chili oil's pungent molecules favor oil and need it to release their heat fully.

But ``like dissolves like'' is a mnemonic, not a law, with limits worth knowing. Ethanol mixes with water in every proportion because it has a water-loving hydroxyl end and a short carbon chain. Longer-chained butanol dissolves only partly; increasing carbon-chain length makes the molecule more ``oily'' and less welcome. Alcohols beyond six carbons are almost excluded like oil. Water affinity is not a label but a contest between the proportions of water-loving and water-avoiding regions. Chapter 25 details saturation and temperature and pressure effects.

Now oil. Cooking-oil molecules consist mainly of long hydrocarbon chains. Carbon and hydrogen have almost identical electronegativities, leaving no pronounced poles along the chain: an ``electrically neutral nice guy.'' What happens in water? Neither substance actually pushes the other away. Oil and water still attract through ordinary dispersion, just weakly. The obstacle lies with water: an oil molecule interrupts the hydrogen-bond net like an object unable to hold hands. Surrounding water forms a restrictive ``cage,'' greatly reducing its possible arrangements. Dispersing oil forces countless water molecules to abandon free rearrangement and line up, a statistically unlikely state to arise spontaneously. Chapter 28 will formalize this intuition as entropy. Energy and probability therefore vote for water clustering together and excluding oil. Droplets merge and layers form, minimizing trapped water molecules.

Remember: \textbf{the hydrophobic effect is not a special ``hydrophobic bond'' between oil molecules}. Oil has only ordinary dispersion. The driving force is the water network's ``self-preservation.'' Oil is not pulled together by something special; water's social network ``shows it out.''

\begin{shiyan}
\textbf{Materials}: two identical small cups, equal quantities of water and cooking oil (about half a cup each), two thermometers (or fingers), and a sunny windowsill.

\textbf{Procedure}: set the cups side by side in sunlight for 30 minutes, measuring every 10 minutes or comparing cup-wall temperatures with a light fingertip touch. Bring them back, leave them together, and compare which cools first.

\textbf{What to observe}: which warms faster, which cools faster, and by how much?

\textbf{Safety note}: work at ordinary temperatures without heating equipment. Do not use glass thermometers for hot or boiling water. Dispose of the oil with food waste, not down a drain.

You will see oil heat and cool faster than water. You have felt the hydrogen-bond network's ``thermal inertia.'' Mild coastal climates and your body's ability to remain at $37\,^{\circ}\mathrm{C}$ are larger-scale versions.
\end{shiyan}

\section{Only Now Do the Symbols Arrive}

Waiting this long to write a formula was deliberate: meaning first, notation second.

Water's formula, \chem{H_2O}, means one oxygen and two hydrogens. To show polarity, chemists place little Greek letters above its ends:
\[\chem{H^{\delta+}-O^{\delta-}-H^{\delta+}}\]
$\delta$, pronounced ``delta,'' means ``a little'': $\delta+$ slightly positive, $\delta-$ slightly negative. These are \textbf{not} complete ionic charges; electrons have not been stolen, only shifted toward oxygen while shared. This differs from Chapter 18's \chem{Na^+} and \chem{Cl^-}, where whole electrons change ownership.

Hydrogen bonds use three dots between molecules:
\[\chem{O-H\cdots O}\]
Solid lines are internal covalent bonds; dots are intermolecular hydrogen bonds. For a structural diagram, a tetrahedral grid with each oxygen holding four neighbors suffices. Remember that it is only a representation: real hydrogen bonds continually break and reconnect, and molecules have neither colors nor little sticks.

\section{How Scientists Know Ice Has Open Spaces}

\begin{zhengju}
Everyone has seen ice float, but the reason provoked years of debate. The simplest guess was trapped air, yet pure ice floats too. Some said solids were simply lighter, although almost every other substance disagreed. A real answer required \textbf{seeing} ice's atomic arrangement. But atoms were too small for any microscope. What could be done?

In 1912, Laue realized that crystals' regular atomic spacings resemble X-ray wavelengths. A crystal should act on X-rays like a grating on light, producing \textbf{diffraction}, an orderly pattern on a photographic plate. Working backward reveals atomic positions. The Braggs soon turned this into a practical ``atomic camera.''

Point it at ice and the mystery resolves: oxygen atoms form a tetrahedral network containing large cavities. Ice really is open. A trickier question remained: hydrogen is too light for X-rays to see readily, so where exactly does it sit between oxygens? Later hydrogen-sensitive neutron diffraction supplied the last piece: hydrogen lies along the oxygen--oxygen line, closer to one end, exactly our \chem{O-H\cdots O} picture. In the 1930s, Pauling organized these clues into the hydrogen-bond concept and used it to explain another strange observation. Hydrogen positions in ordinary ice are somewhat disordered, retaining disorder even at absolute zero. Counting hydrogen's two possible placements between \chem{O\cdots O} gave quantitative results almost exactly matching experiments.

The story continues. Liquid water is harder to picture than ice: its flickering rearrangements evade still frames. Even today, the exact shape of its hydrogen-bond network is an active research topic, with different experimental pictures still competing. This is real science: ice's floating is settled, while network details advance through debate.
\end{zhengju}

\section{Even This Network Has Limits}

``Polar molecules plus a hydrogen-bond net'' is useful, but not the whole truth about water.

First, \textbf{there is more than one ice}. Familiar ice forms at ordinary pressure. Other temperatures and pressures create entirely different structures; scientists have confirmed more than a dozen crystalline forms. Some are denser than liquid water and would sink. ``Ice always floats'' describes ordinary-pressure ice, not water's universal nature.

Second, \textbf{hydrogen bonding has blurred boundaries}. Its strength spans a continuum between chemical bonds and ordinary intermolecular interactions, and textbooks draw different lines. A region on a strength spectrum is closer to reality than a black-and-white category.

Third, \textbf{do not follow marketing language astray}. ``Small-cluster water,'' ``activated water,'' and ``water memory'' occasionally advertise healthier treated water. The network does change, but unimaginably fast: any structure is transformed on trillionth-of-a-second timescales, unable to store information like a recorded disc. Water remembering is like a bonfire remembering each log's shape. Chapter 25 also shows that swallowed water immediately joins body water, without its prior ``treatments'' following along.

\begin{wuqu}
``Oil and water do not mix because special hydrophobic bonds bind oil molecules tightly.'' This reverses cause and effect. Disperse oil into air as a mist and droplets still merge, although nothing there is ``hydrophobic.'' They merge because surface molecules are pulled inward (Chapter 22's surface tension). In water, too, oil has no special glue. Water's effort to restore its network drives oil aside and together. Test an explanation by asking whose ``push'' it invokes. Assigning it to oil usually signals a misconception; assigning it to water's network gives the modern scientific answer.
\end{wuqu}

\section{Back to the Ice Water}

Now settle the three opening mysteries.

Floating ice: cavities in its tetrahedral scaffold lower density to about nine-tenths of liquid water's, leaving a tenth above the surface---literally an iceberg's tip.

Droplets outside: countless water molecules already drift through air. Hitting the cold glass slows them until they cannot escape the network's capture. They gather and condense into visible drops, recruited from air rather than leaked from inside.

Unchanged level after melting: floating ice displaces exactly the water it becomes. The melted water fills precisely its former underwater ``hole,'' leaving the level almost unchanged. (Sea ice likewise does not raise sea level when melting; land ice sheets sliding into the sea cause the rise.)

Add our everyday questions: warm-water sponging reduces fever through evaporation's generous \ud{2400}{J} per gram, while coastal mildness comes from seawater's large thermal inertia. In an ordinary glass of ice water, the hydrogen-bond network performs every act.

\section{What Lies Beyond the Network?}

We have seen \textbf{one kind of force, organized differently, produce phenomena across the scales of liquids, solids, climate, and body temperature}. Water is special not because one property is exceptionally strong, but because its tetrahedral net is ``just right'': strong enough to keep molecules liquid, weak enough to rearrange constantly, dissolving salt and sugar while excluding oil.

The story continues in two directions. Downstream: what happens when salt, sugar, gases, and other substances enter the net? How do dissolved particles arrange themselves, conduct, and compete across semipermeable membranes? Chapter 25, ``What Happens in Solutions,'' covers this. Upstream: put a molecule with one water-loving and one oil-loving end into water, and what formation do the friendship rules demand? Astonishingly, they spontaneously surround hollow spheres, hiding oil-loving ends in the walls. Billions of years ago, just such a self-assembled underwater bubble first separated inside from outside, becoming a cell boundary. Chapter 51's membranes return to this network, revealing life's boundary as another consequence of water's pickiness.

\begin{tiaozhan}
One water anomaly remains: liquid water is densest at $4\,^{\circ}\mathrm{C}$, becoming less dense whether heated or cooled. Why? Two tendencies compete. Cooling weakens thermal motion, letting hydrogen bonds arrange neighbors at standard distances. Greater order and closer spacing increase density. Simultaneously, more regions form ice-like tetrahedral clusters whose cavities expand the volume. Near $4\,^{\circ}\mathrm{C}$, the second tendency catches the first and density peaks. Further cooling lets cavities dominate; approaching freezing loosens water. Thus $0\,^{\circ}\mathrm{C}$ water is already lighter than $4\,^{\circ}\mathrm{C}$ water, and winter lake bottoms remain near $4\,^{\circ}\mathrm{C}$, sheltering aquatic life. For numbers, look up densities of about \ud{1.000}{g/cm^3} at $4\,^{\circ}\mathrm{C}$ and \ud{0.9998}{g/cm^3} at $0\,^{\circ}\mathrm{C}$. A tiny difference, yet enough to drive whole-lake turnover before winter and carry oxygen to the bottom. Small numbers, big ecology.
\end{tiaozhan}

\section*{Try It}

\begin{sikaoti}
\item Draw a water molecule with $\delta+$ and $\delta-$ ends. Add three neighbors, using solid lines for covalent bonds and three dots for hydrogen bonds, so the central molecule holds exactly four ``hands.'' Note that this is a thinking aid; real hydrogen bonds continually break and reconnect.
\item Ammonia (\chem{NH_3}) also hydrogen-bonds, but each molecule holds fewer neighbors on average than water. Predict its boiling point relative to water and methane, then check data.
\item A river freezes at its surface but not its bottom in winter. Draw a cross section marking ice, the $4\,^{\circ}\mathrm{C}$ water layer, and fish. Explain structurally, and say what would happen if ice were heavier than water.
\item Use the network's self-preservation to explain step by step why shaken oil and water separate after standing a few minutes. Avoid ``oil hates water'' or ``hydrophobic bonds''; start with water's energy and arrangements.
\item Evaporating 1 gram of sweat removes about \ud{2400}{J}, while cooling 1 gram of water by $1\,^{\circ}\mathrm{C}$ releases only \ud{4.2}{J}. Estimate how much evaporating 10 grams could theoretically cool a 60-kilogram person, treating the body entirely as water. Why is the actual effect smaller?
\item Someone claims a brand's treated water has smaller molecular clusters that cells absorb better. Write three sentences refuting this using the network's rearrangement speed.
\end{sikaoti}
